% % Methane (CH\textsubscript{4}) is a potent greenhouse gas with approximately
% % 28 times the global warming potential of CO\textsubscript{2} over 100 years,
% % and agriculture is a major anthropogenic source. Eddy Covariance measurements
% % capture ecosystem-scale emissions from managed grasslands, but these records
% % are often incomplete and dominated by episodic emission peaks, making methane
% % flux difficult to reconstruct and predict.

% Methane (CH\textsubscript{4}) is a potent greenhouse gas, with approximately 28
% times the global warming potential of CO\textsubscript{2} over 100 years. UK
% agriculture accounts for almost half (49\%) of national methane emissions 
% approximately 27 MtCO\textsubscript{2}e in 2024 with ruminant enteric fermentation
% forming the principal agricultural source. Eddy Covariance measurements provide
% ecosystem-scale observations from managed grasslands, but the resulting records are
% often incomplete and dominated by episodic emission peaks. Gap-filling supports the
% construction of continuous historical flux records, while forecasting provides a
% basis for anticipating future emissions and informing management decisions.

Methane (CH\textsubscript{4}) has approximately 28 times the 100-year global warming
potential of CO\textsubscript{2}. UK agriculture contributes almost half (49\%) of 
national methane emissions. Eddy 
Covariance (EC) provides ecosystem-scale observations of methane exchange over managed 
grasslands. 
% enabling data-driven modelling of environmental and management influences. 
However, EC records are often incomplete and characterised by episodic emission peaks. 
Gap-filling can reconstruct continuous historical flux records, while forecasting can 
anticipate future emissions and inform agricultural management.

This dissertation develops a three-stage machine-learning pipeline for
gap-filling, forecasting, and scenario projection using data from three
catchments at the North Wyke Farm Platform. Statistical, tree-based, neural,
and tabular pretrained foundation models were benchmarked using held-out and temporal
validation, supported by uncertainty quantification and permutation
importance. Future climate trajectories and emission scenarios were then combined with  
alternative farm management conditions to explore possible changes in methane flux.

% TabICLv2 achieved the strongest gap-filling performance, increasing combined 
% hourly $R^2$ from 0.055 for marginal distribution sampling to 0.515. TabPFN 
% led forecasting with MASE$=0.6908$, although performance deteriorated from 
% 0.524 on ordinary days to 3.285 during methane spikes. 

% The pretrained tabular models provided the strongest overall performance for 
% both gap-filling and forecasting. Among the models tested, traditional tree-based 
% models generally outperformed neural networks with the exception of tabular pretrained models. 

Overall performance follows a consistent pattern of tabular pretrained models performing best,
followed by tree-based models, and locally trained neural networks generally underperforming.
Nevertheless, sudden methane peaks remained difficult to anticipate for all models. 
Among the tested scenarios, livestock density produced the largest change in predicted annual 
CH\textsubscript{4} flux, followed by grazing duration, whereas varying fertiliser 
application rates produced little to no change in the modelled response.

Overall, pretrained foundation models for tabular data performed best across the predictive tasks,
while livestock-associated variables formed the strongest recurring influence
at the livestock-dominated catchments. However, severe target missingness,
emission spikes, and the absence of future ground truth limit 
the operational utility of these models.
